\section{Architectures}
For this project I have used three workstation, my personal one, the workstation provided by the University and Google Colab.

\subsection{Personal Workstation}
\textbf{HP Envy 17-ae103nl}

\begin{table}[h!]
\renewcommand{\arraystretch}{1.2}
\begin{tabularx}{\textwidth}{lX}
\multicolumn{2}{l}{\textbf{CPU}} \\ \hline
Model & Intel Core i7-6500U @ 2.50 GHz \\
Architecture & x86\_64 \\
Physical cores & 2 \\
Logical threads & 4 (Hyper-Threading) \\[6pt]

\multicolumn{2}{l}{\textbf{GPU}} \\ \hline
Model & NVIDIA GeForce GTX 950M \\
CUDA Compute Capability & 5.0 \\
Driver version & 572.16 \\
CUDA Version & 12.8 \\[6pt]

\multicolumn{2}{l}{\textbf{Software and Setup}} \\ \hline
Visual Studio & Visual Studio Community 2022  \\
                & (C++ Desktop workload) \\ 
CUDA Toolkit & Installed manually \\
WSL2 & Enabled \\
NVIDIA Drivers & Version 572.16 (for WSL2) \\
\end{tabularx}
\end{table}


\subsection{University Workstation}
\textbf{Room 210}

\begin{table}[h!]
\renewcommand{\arraystretch}{1.2}
\begin{tabularx}{\textwidth}{lX}

\multicolumn{2}{l}{\textbf{CPU}} \\ \hline
Model & Intel Core i9-12900K (12th Gen) \\
Architecture & x86\_64 \\
CPU(s) & 24 \\
Threads per core & 2 \\
Cores per socket & 16 \\
Socket(s) & 1 \\
Max frequency & 4.02 GHz \\
Profiling tools & Intel VTune Profiler, Intel Advisor \\[6pt]

\end{tabularx}
\end{table}

To use the Intel oneAPI tools, the environment was loaded with:
\begin{verbatim}
        source /opt/intel/oneapi/setvars.sh
\end{verbatim}


\subsection{Google Colab}

\begin{table}[h!]
\renewcommand{\arraystretch}{1.2}
\begin{tabularx}{\textwidth}{lX}

\multicolumn{2}{l}{\textbf{CPU}} \\ \hline
Model & Intel(R) Xeon(R) CPU @ 2.20GHz \\
Architecture & x86\_64 \\
Physical cores & 1 \\
Logical threads & 2 (Hyper-Threading) \\
Socket(s) & 1 \\[6pt]

\multicolumn{2}{l}{\textbf{GPU}} \\ \hline
Model & NVIDIA Tesla T4 \\
CUDA Compute Capability & 7.5 \\
Driver version & 550.54.15 \\
CUDA Version & 12.8 \\[6pt]

\end{tabularx}
\end{table}
